\chapter*{Appendix A\\ Framework configurations}
\label{appendix}
In this section we are going to report the followed procedures for installing and setting up the frameworks used in this thesis work.
\section{Apache Oozie}

The version used for this project is the 4.2.0 (stable version - June 2015). The procedure for the installation is not trivial, as you need to build a working distro from source code.
Prerequisites:
\begin{itemize}
\item Java 1.6+
\item Apache Hadoop
\item Apache Maven
\item ext-2.2
\end{itemize}
Java 1.8, hadoop 2.7.0 and maven 3.3.3 were used during the installation procedure.\\
It is assumed also that hadoop is running using the default ports setting
\subsection{Installation\protect\footnote{At the time of installation I came across a bug due to a repository out of service. Refer to \cite{bib3} for its solution}}
\begin{lstlisting}[language=sh]
#donwload the tarball from a mirror
cd ~
wget http://apache.uib.no/oozie/4.2.0/oozie-4.2.0.tar.gz 
tar -xvf oozie-4.2.0.tar.gz

#build a distro according to your java version and hadoop version
cd oozie-4.2.0/bin
./mkdistro -DtargetJavaVersion=1.8 -P hadoop-2 -Dhadoop.version=2.7.0 

#configure the distro
cd distro/target/oozie-4.2.0
mkdir libext
mv /path/to/ext-2.2.zip libext
bin/oozie-setup.sh sharelib create -fs hdfs://localhost:9000
bin/oozie-setup.sh prepare-war
bin/ooziedb.sh create -sqlfile oozie.sql -run
bin/oozied.sh start
\end{lstlisting}

In order to set up the integration with hadoop, the following properties need to be defined:

\begin{lstlisting}[language=XML]
<!-- oozie-site.xml -->
<property>
    <name>oozie.service.HadoopAccessorService.hadoop.configurations</name>
	<value>*=/path/to/hadoop-2.7.0/etc/hadoop</value>
</property>
<property>
     <name>oozie.service.WorkflowAppService.system.libpath</name>
     <value>hdfs:///user/${user.name}/share/lib</value>
</property>

<!-- hadoop core-site.xml -->
<property>
      <name>hadoop.proxyuser.myhadoopusername.hosts</name>
      <value>*</value>
</property>
<property>
      <name>hadoop.proxyuser.myhadoopusername.groups</name>es
      <value>*</value>
</property>

\end{lstlisting}

Once oozie is started, is it possible to control the job executions via web interface at $localhost:11000$


 